\section{Norms and Distances}

\begin{res}{Lemma}{4.1}
Any two nonnegative numbers $a,b\in\mathbb R_{\ge0}$ satisfy \textbf{Young's inequality}
\begin{equation}ab\le\frac{a^p}{p}+\frac{b^q}{q}.\tag{4.4}\end{equation}
\end{res}

\begin{res}{Lemma}{4.2}
Any two vectors $\boldsymbol v,\boldsymbol w\in\mathbb F^n$ satisfy \textbf{H\"older's inequality}
\begin{equation}\sum_{i=1}^n\left|v^i\right|\left|w^i\right|\le\|\boldsymbol v\|_p\|\boldsymbol w\|_q\tag{4.5}\end{equation}
\end{res}

\begin{res}{Proposition}{4.3}
Any two vectors $\boldsymbol v,\boldsymbol w\in\mathbb F^n$ satisfy \textbf{Minkowski's inequality}
\begin{equation}\|\boldsymbol v+\boldsymbol w\|_p\le\|\boldsymbol v\|_p+\|\boldsymbol w\|_p.\tag{4.6}\end{equation}
\end{res}

\begin{res}{Proposition}{4.4}
A normed vector space $(V,\|\cdot\|)$ becomes a metric space using the definition $d_V:V\times V\to\mathbb R:(v,w)\mapsto\|v-w\|$.
\end{res}

\begin{res}{Proposition}{4.5}
The norm itself is a continuous function $V\to\mathbb R$.
\end{res}

\begin{res}{Proposition}{4.6}
Vector addition is a continuous map from $V\times V$ to $V$, and scalar multiplication is a continuous map from $\mathbb F\times V$ to $V$.
\end{res}

\begin{res}{Proposition}{4.7}
Two norms $\|v\|_a$ and $\|v\|_b$ on $V$ are equivalent if and only if there exist constants $c,C>0$ such that
\begin{equation}\forall v\in V:c\,\|v\|_a\le\|v\|_b\le C\,\|v\|_a.\tag{4.10}\end{equation}
\end{res}

\begin{res}{Proposition}{4.8}
Let $V$ be a finite-dimensional vector space. Any two norms on $V$ are equivalent.
\end{res}

\begin{res}{Corollary}{4.9}
A consequence is that on any finite-dimensional vector space $V$, any norm is continuous with respect to the metric generated by any other norm.
\end{res}

\begin{res}{Proposition}{4.10}
For a metric space $(X,d_X)$, any sequence $(x_n\in X)_{n\in\mathbb N_0}$ that converges to a limit $x$ is a Cauchy sequence.
\end{res}

\begin{res}{Proposition}{4.11}
A normed space $(V,\|\cdot\|)$ is metric complete (and thus a Banach space) if and only if every absolutely convergent series converges.
\end{res}

\begin{res}{Proposition}{4.12}
For a linear map $\hat A\in\operatorname{Hom}(V,W)$, the following are equivalent:
\begin{enumerate}
\item $\hat A$ is bounded
\item $\hat A$ is continuous at the origin (or some other point).
\item $\hat A$ is continuous everywhere, and is in fact uniformly continuous.
\end{enumerate}
\end{res}

\begin{res}{Proposition}{4.13}
Let $(V,\|\cdot\|_V)$ and $(W,\|\cdot\|_W)$ be normed vector spaces. If $V$ is finite-dimensional, then $\mathcal B(V,W)=\operatorname{Hom}(V,W)$.
\end{res}

\begin{res}{Proposition}{4.14}
$(V,\|\cdot\|_V)$ be a normed vector space and $(W,\|\cdot\|_W)$ be a Banach space. The space of bounded linear maps $\mathcal B(V,W)$ together with the induced norm $\|\cdot\|_{V\to W}$ is a Banach space.
\end{res}

\begin{res}{Proposition}{4.15}
Let $(V,\|\cdot\|)$ be a normed vector space. The associated operator norm $\|\cdot\|_{V\to V}$ on $\operatorname{End}(V)$ is submultiplicative.
\end{res}

\begin{res}{Proposition}{4.16}
The norms $\|\mathsf A\|_p=\|\mathsf A\|_{p\to p}$ that are induced by the vector $p$-norm can be defined for any matrix size, and form a consistent family of matrix norms.
\end{res}

\begin{res}{Proposition}{4.17}
The Frobenius norm $\mathsf A_{\mathrm F}$ is consistent, and thus submultiplative and subordinate with respect to the vector 2-norm, to which it reduces for column matrices in $\mathbb F^{n\times1}$.
\end{res}

\begin{res}{Proposition}{4.18}
Let $\|\cdot\|$ be a submultiplicative norm on $\operatorname{End}(V)$ for a finite-dimensional vector space $V$. For any $\hat A\in\operatorname{End}(V)$, we find $\|\hat A\|\ge\rho_{\hat A}$.
\end{res}

\begin{res}{Proposition}{4.19}
Let $\|\cdot\|$ be a submultiplicative norm on $\operatorname{End}(V)$ for a finite-dimensional vector space $V$. For any $\hat A\in\operatorname{End}(V)$, we have
\begin{equation}\lim_{n\to\infty}\|\hat A^n\|^{1/n}=\rho_{\hat A}.\tag{4.24}\end{equation}
This result is known as the \textbf{Gelfand formula}.
\end{res}
